% Faithful transcription — original German. Content and notation follow the 1868 printing;
% only markup is standardized (PLAN.md §4.4). This fixture reproduces the opening of the
% Plan der Untersuchung; it is a short excerpt for demonstration, not the full memoir.
\documentclass{article}
\usepackage{readmasters}
\begin{document}

\origpage{133}
\section*{Plan der Untersuchung}

Bekanntlich setzt die Geometrie sowohl den Begriff des Raumes, als die ersten
Grundbegriffe für die Constructionen im Raume als etwas Gegebenes voraus. Sie giebt
von ihnen nur Nominaldefinitionen, während die wesentlichen Bestimmungen in Form von
Axiomen auftreten.

Das Verhältniss dieser Voraussetzungen bleibt dabei im Dunkeln; man sieht weder ein,
ob und in wie weit ihre Verbindung nothwendig, noch a~priori, ob sie möglich ist.

\origpage{134}
Diese Dunkelheit wurde auch von Euklid bis auf Legendre, um den berühmtesten neueren
Bearbeiter der Geometrie zu nennen, weder von den Mathematikern, noch von den
Philosophen, welche sich damit beschäftigten, gehoben. Es hatte dies seinen Grund
wohl darin, dass der allgemeine Begriff mehrfach ausgedehnter Grössen, unter welchem
die Raumgrössen enthalten sind, ganz unbearbeitet blieb.

Ich habe mir daher zunächst die Aufgabe gestellt, den Begriff einer mehrfach
ausgedehnten Grösse aus allgemeinen Grössenbegriffen zu construiren.

\end{document}
